% =============================================================
%  Chapter 4 — nGibbs / ML Thermodynamics
%  "New Frontiers in Thermodynamically Informed Geodynamic
%   Modeling: a General Deep Learning Accelerated Free Energy
%   Minimizer for Multicomponent Silicate Systems"
%
%  Relocated from Chapter5/chapter5.tex on 2026-09-01: DustMAP (the
%  Cassini CDA/INMS pipeline chapter that previously occupied the
%  Chapter4 slot) has been pulled from the thesis and shelved at
%  Manuscripts/DustMAP/ as a standalone manuscript-in-progress instead.
%  This nGibbs chapter moves up into Chapter4 as a result; Chapter5 is
%  now the Draaikraal (Bushveld margin) field-study chapter. The
%  \label{ch:ngibbs} below is unchanged, so ChapterA/chapterA.tex's
%  \ref{ch:ngibbs} in its Conclusions still resolves correctly.
%
%  This stub will be replaced with the full chapter text.
% =============================================================

\chapter{New Frontiers in Thermodynamically Informed Geodynamic Modeling: a General Deep Learning Accelerated Free Energy Minimizer for Multicomponent Silicate Systems}
\label{ch:ngibbs}

\begin{quote}
  \itshape
  \todo{One-paragraph chapter summary.}
\end{quote}

\section{Introduction}
\label{sec:ch4-intro}
\todo{Motivate the need for fast thermodynamic forward models in geodynamics.
Review current tools (Perple\_X, THERMOCALC, MAGEMin) and their computational
bottlenecks. State the deep learning acceleration approach.}

\section{Background}
\label{sec:ch4-background}

\subsection{Gibbs Free Energy Minimization in Petrology}
\todo{}

\subsection{Machine Learning Surrogates for Physical Models}
\todo{}

\section{Model Architecture}
\label{sec:ch4-architecture}

\subsection{Training Data Generation}
\todo{Describe the strategy for generating phase equilibrium training data
from classical minimizers across relevant $P$--$T$--$X$ space.}

\subsection{Neural Network Design}
\todo{Describe the architecture: input/output spaces, layers, activation
functions, loss function, training procedure.}

\subsection{Uncertainty Quantification}
\todo{}

\section{Validation}
\label{sec:ch4-validation}
\todo{Benchmark speed and accuracy against classical minimizers on held-out
$P$--$T$--$X$ points and on known phase diagrams.}

\section{Application to Geodynamic Modeling}
\label{sec:ch4-results}
\todo{Demonstrate integration into a geodynamic simulation. Show a
representative case study.}

\section{Discussion}
\label{sec:ch4-discussion}
\todo{Discuss accuracy/speed tradeoffs, generalizability to other
bulk compositions, and extension to systems relevant to the Bushveld
Complex (foreshadow Appendix~\ref{app:bushveld} and the Conclusion).}

\section{Conclusions}
\label{sec:ch4-conclusions}
\todo{Summarize the ML minimizer performance and implications for
future thermodynamically coupled geodynamic models.}
